% 原书例目录：印刷xxiii–xxix／PDF物理24–30；原编号与原页保留.
\clearpage\phantomsection
\section*{例目录（原书页码）}\phantomsection\label{front:examples}
\addcontentsline{toc}{section}{例目录（原书页码）}
\begin{longtable}{@{}p{.82\textwidth}r@{}}
例 1.1.3（事件运算） & \hyperref[ex:1.1.3]{2} \\
例 1.2.2（$\sigma$代数—I） & \hyperref[ex:1.2.2]{5} \\
例 1.2.3（$\sigma$代数—II） & \hyperref[ex:1.2.3]{5} \\
例 1.2.5（定义概率—I） & \hyperref[ex:1.2.5]{6} \\
例 1.2.7（定义概率—II） & \hyperref[ex:1.2.7]{7} \\
例 1.2.10（Bonferroni不等式） & \hyperref[ex:1.2.10]{9} \\
例 1.2.12（彩票—I） & \hyperref[ex:1.2.12]{11} \\
例 1.2.13（锦标赛） & \hyperref[ex:1.2.13]{11} \\
例 1.2.15（彩票—II） & \hyperref[ex:1.2.15]{11} \\
例 1.2.18（扑克牌） & \hyperref[ex:1.2.18]{13} \\
例 1.2.19（有放回抽样） & \hyperref[ex:1.2.19]{14} \\
例 1.2.20（计算平均值） & \hyperref[ex:1.2.20]{15} \\
例 1.3.1（四张A） & \hyperref[ex:1.3.1]{16} \\
例 1.3.3（例1.3.1续） & \hyperref[ex:1.3.3]{17} \\
例 1.3.4（三个囚犯） & \hyperref[ex:1.3.4]{18} \\
例 1.3.6（编码） & \hyperref[ex:1.3.6]{19} \\
例 1.3.8（de Méré骑士） & \hyperref[ex:1.3.8]{20} \\
例 1.3.10（掷两枚骰子） & \hyperref[ex:1.3.10]{21} \\
例 1.3.11（信件） & \hyperref[ex:1.3.11]{21} \\
例 1.3.13（抛掷硬币三次—I） & \hyperref[ex:1.3.13]{22} \\
例 1.4.2（随机变量） & \hyperref[ex:1.4.2]{23} \\
例 1.4.3（抛掷硬币三次—II） & \hyperref[ex:1.4.3]{23} \\
例 1.4.4（随机变量的分布） & \hyperref[ex:1.4.4]{24} \\
例 1.5.2（抛掷三枚硬币） & \hyperref[ex:1.5.2]{24} \\
例 1.5.4（抛掷直到正面出现） & \hyperref[ex:1.5.4]{25} \\
例 1.5.5（连续的累积分布函数） & \hyperref[ex:1.5.5]{26} \\
例 1.5.6（有跳跃的累积分布函数） & \hyperref[ex:1.5.6]{27} \\
例 1.5.9（同分布的随机变量） & \hyperref[ex:1.5.9]{28} \\
例 1.6.2（几何分布概率） & \hyperref[ex:1.6.2]{28} \\
例 1.6.4（Logistic分布概率） & \hyperref[ex:1.6.4]{30} \\
\addlinespace[.5em]
例 2.1.1（二项分布的变换） & \hyperref[ex:2.1.1]{39} \\
例 2.1.2（均匀分布的变换） & \hyperref[ex:2.1.2]{39} \\
例 2.1.4（均匀分布与指数分布的关系—I） & \hyperref[ex:2.1.4]{41} \\
例 2.1.6（逆伽马概率密度函数） & \hyperref[ex:2.1.6]{42} \\
例 2.1.7（平方变换） & \hyperref[ex:2.1.7]{42} \\
例 2.1.9（正态分布与卡方分布的关系） & \hyperref[ex:2.1.9]{43} \\
例 2.2.2（指数分布的均值） & \hyperref[ex:2.2.2]{45} \\
例 2.2.3（二项分布的均值） & \hyperref[ex:2.2.3]{45} \\
例 2.2.4（Cauchy分布的均值） & \hyperref[ex:2.2.4]{46} \\
例 2.2.6（使距离最小） & \hyperref[ex:2.2.6]{47} \\
例 2.2.7（均匀分布与指数分布的关系—II） & \hyperref[ex:2.2.7]{48} \\
例 2.3.3（指数分布的方差） & \hyperref[ex:2.3.3]{49} \\
例 2.3.5（二项分布的方差） & \hyperref[ex:2.3.5]{50} \\
例 2.3.8（伽马分布的矩母函数） & \hyperref[ex:2.3.8]{51} \\
例 2.3.9（二项分布的矩母函数） & \hyperref[ex:2.3.9]{52} \\
例 2.3.10（不能唯一确定分布的矩） & \hyperref[ex:2.3.10]{53} \\
例 2.3.13（Poisson近似） & \hyperref[ex:2.3.13]{54} \\
例 2.4.5（交换积分与微分—I） & \hyperref[ex:2.4.5]{58} \\
例 2.4.6（交换积分与微分—II） & \hyperref[ex:2.4.6]{59} \\
例 2.4.7（交换求和与微分） & \hyperref[ex:2.4.7]{60} \\
例 2.4.9（习题2.4.7续）\footnote{译注：原例目录写“Exercise 2.4.7”，但2.4.7是例的编号，原目录用词保留.} & \hyperref[ex:2.4.9]{61} \\
\addlinespace[.5em]
例 3.2.1（验收抽样） & \hyperref[ex:3.2.1]{72} \\
例 3.2.3（掷骰子的概率） & \hyperref[ex:3.2.3]{74} \\
例 3.2.4（等待时间） & \hyperref[ex:3.2.4]{76} \\
例 3.2.5（Poisson近似） & \hyperref[ex:3.2.5]{77} \\
例 3.2.6（逆二项抽样） & \hyperref[ex:3.2.6]{79} \\
例 3.2.7（失效时间） & \hyperref[ex:3.2.7]{80} \\
例 3.3.1（伽马分布与Poisson分布的关系） & \hyperref[ex:3.3.1]{82} \\
例 3.3.2（正态近似） & \hyperref[ex:3.3.2]{86} \\
例 3.4.1（二项指数族） & \hyperref[ex:3.4.1]{92} \\
例 3.4.3（二项分布的均值与方差） & \hyperref[ex:3.4.3]{93} \\
例 3.4.4（正态指数族） & \hyperref[ex:3.4.4]{93} \\
例 3.4.6（例3.4.4续） & \hyperref[ex:3.4.6]{94} \\
例 3.4.8（一个弯曲指数族） & \hyperref[ex:3.4.8]{95} \\
例 3.4.9（正态近似） & \hyperref[ex:3.4.9]{95} \\
例 3.5.3（指数分布的位置族） & \hyperref[ex:3.5.3]{97} \\
例 3.6.2（Chebychev不等式的例示） & \hyperref[ex:3.6.2]{101} \\
例 3.6.3（一个正态概率不等式） & \hyperref[ex:3.6.3]{101} \\
例 3.6.6（正态分布的高阶矩） & \hyperref[ex:3.6.6]{103} \\
例 3.6.9（Poisson分布的高阶矩） & \hyperref[ex:3.6.9]{104} \\
\addlinespace[.5em]
例 4.1.2（掷骰子的样本空间） & \hyperref[ex:4.1.2]{114} \\
例 4.1.4（例4.1.2续） & \hyperref[ex:4.1.4]{116} \\
例 4.1.5（掷骰子的联合概率质量函数） & \hyperref[ex:4.1.5]{116} \\
例 4.1.7（掷骰子的边缘概率质量函数） & \hyperref[ex:4.1.7]{117} \\
例 4.1.8（掷骰子的概率） & \hyperref[ex:4.1.8]{118} \\
例 4.1.9（边缘分布相同、联合概率质量函数不同） & \hyperref[ex:4.1.9]{118} \\
例 4.1.11（计算联合概率—I） & \hyperref[ex:4.1.11]{119} \\
例 4.1.12（计算联合概率—II） & \hyperref[ex:4.1.12]{120} \\
例 4.2.2（计算条件概率） & \hyperref[ex:4.2.2]{122} \\
例 4.2.4（计算条件概率密度函数） & \hyperref[ex:4.2.4]{123} \\
例 4.2.6（检验独立性—I） & \hyperref[ex:4.2.6]{125} \\
例 4.2.8（检验独立性—II） & \hyperref[ex:4.2.8]{126} \\
例 4.2.9（联合概率模型） & \hyperref[ex:4.2.9]{126} \\
例 4.2.11（独立变量的期望） & \hyperref[ex:4.2.11]{127} \\
例 4.2.13（正态变量之和的矩母函数） & \hyperref[ex:4.2.13]{128} \\
例 4.3.1（Poisson变量之和的分布） & \hyperref[ex:4.3.1]{129} \\
例 4.3.3（贝塔变量之积的分布） & \hyperref[ex:4.3.3]{130} \\
例 4.3.4（正态变量的和与差） & \hyperref[ex:4.3.4]{131} \\
例 4.3.6（正态变量之比的分布） & \hyperref[ex:4.3.6]{132} \\
例 4.4.1（二项–Poisson层次模型） & \hyperref[ex:4.4.1]{133} \\
例 4.4.2（例4.4.1续） & \hyperref[ex:4.4.2]{134} \\
例 4.4.5（例4.4.1的推广） & \hyperref[ex:4.4.5]{135} \\
例 4.4.6（贝塔–二项层次模型） & \hyperref[ex:4.4.6]{137} \\
例 4.4.8（例4.4.6续） & \hyperref[ex:4.4.8]{138} \\
例 4.5.4（相关性—I） & \hyperref[ex:4.5.4]{139} \\
例 4.5.8（相关性—II） & \hyperref[ex:4.5.8]{141} \\
例 4.5.9（相关性—III） & \hyperref[ex:4.5.9]{143} \\
例 4.6.1（多元概率密度函数） & \hyperref[ex:4.6.1]{146} \\
例 4.6.3（多元概率质量函数） & \hyperref[ex:4.6.3]{148} \\
例 4.6.8（伽马变量之和的矩母函数） & \hyperref[ex:4.6.8]{150} \\
例 4.6.13（多元变量替换） & \hyperref[ex:4.6.13]{152} \\
例 4.7.4（协方差不等式） & \hyperref[ex:4.7.4]{153} \\
例 4.7.8（一个关于均值的不等式） & \hyperref[ex:4.7.8]{156} \\
\addlinespace[.5em]
例 5.1.2（样本概率密度函数：指数分布） & \hyperref[ex:5.1.2]{170} \\
例 5.1.3（有限总体模型） & \hyperref[ex:5.1.3]{171} \\
例 5.2.8（均值的分布） & \hyperref[ex:5.2.8]{175} \\
例 5.2.10（Cauchy随机变量之和） & \hyperref[ex:5.2.10]{176} \\
例 5.2.12（Bernoulli随机变量之和） & \hyperref[ex:5.2.12]{177} \\
例 5.3.5（方差比的分布） & \hyperref[ex:5.3.5]{182} \\
例 5.3.7（例5.3.5续） & \hyperref[ex:5.3.7]{183} \\
例 5.4.5（均匀分布次序统计量的概率密度函数） & \hyperref[ex:5.4.5]{187} \\
例 5.4.7（中程数与极差的分布） & \hyperref[ex:5.4.7]{188} \\
例 5.5.3（$S^2$的相合性） & \hyperref[ex:5.5.3]{190} \\
例 5.5.5（$S$的相合性） & \hyperref[ex:5.5.5]{190} \\
例 5.5.7（几乎处处收敛） & \hyperref[ex:5.5.7]{191} \\
例 5.5.8（依概率收敛但不几乎处处收敛） & \hyperref[ex:5.5.8]{191} \\
例 5.5.11（均匀随机变量的最大值） & \hyperref[ex:5.5.11]{192} \\
例 5.5.16（负二项分布的正态近似） & \hyperref[ex:5.5.16]{194} \\
例 5.5.18（使用估计方差的正态近似） & \hyperref[ex:5.5.18]{195} \\
例 5.5.19（估计胜算） & \hyperref[ex:5.5.19]{195} \\
例 5.5.22（例5.5.19续） & \hyperref[ex:5.5.22]{197} \\
例 5.5.23（近似均值与方差） & \hyperref[ex:5.5.23]{197} \\
例 5.5.25（例5.5.23续） & \hyperref[ex:5.5.25]{198} \\
例 5.5.27（比率估计量的矩） & \hyperref[ex:5.5.27]{199} \\
例 5.6.1（指数分布的寿命） & \hyperref[ex:5.6.1]{200} \\
例 5.6.2（例5.6.1续） & \hyperref[ex:5.6.2]{200} \\
例 5.6.3（概率积分变换） & \hyperref[ex:5.6.3]{201} \\
例 5.6.4（Box–Muller，1958） & \hyperref[ex:5.6.4]{202} \\
例 5.6.5（二项随机变量的生成） & \hyperref[ex:5.6.5]{203} \\
例 5.6.6（Poisson方差的分布） & \hyperref[ex:5.6.6]{203} \\
例 5.6.7（贝塔随机变量的生成—I） & \hyperref[ex:5.6.7]{204} \\
例 5.6.9（贝塔随机变量的生成—II） & \hyperref[ex:5.6.9]{207} \\
\addlinespace[.5em]
例 6.2.3（二项分布的充分统计量） & \hyperref[ex:6.2.3]{223} \\
例 6.2.4（正态分布的充分统计量） & \hyperref[ex:6.2.4]{224} \\
例 6.2.5（充分的次序统计量） & \hyperref[ex:6.2.5]{224} \\
例 6.2.7（例6.2.4续） & \hyperref[ex:6.2.7]{225} \\
例 6.2.8（均匀分布的充分统计量） & \hyperref[ex:6.2.8]{226} \\
例 6.2.9（两个参数均未知时正态分布的充分统计量） & \hyperref[ex:6.2.9]{227} \\
例 6.2.12（两个正态充分统计量） & \hyperref[ex:6.2.12]{228} \\
例 6.2.14（正态分布的极小充分统计量） & \hyperref[ex:6.2.14]{229} \\
例 6.2.15（均匀分布的极小充分统计量） & \hyperref[ex:6.2.15]{229} \\
例 6.2.17（均匀分布的辅助统计量） & \hyperref[ex:6.2.17]{230} \\
例 6.2.18（位置族的辅助统计量） & \hyperref[ex:6.2.18]{231} \\
例 6.2.19（尺度族的辅助统计量） & \hyperref[ex:6.2.19]{231} \\
例 6.2.20（辅助精度） & \hyperref[ex:6.2.20]{232} \\
例 6.2.22（二项分布的完备充分统计量） & \hyperref[ex:6.2.22]{232} \\
例 6.2.23（均匀分布的完备充分统计量） & \hyperref[ex:6.2.23]{233} \\
例 6.2.26（使用Basu定理—I） & \hyperref[ex:6.2.26]{235} \\
例 6.2.27（使用Basu定理—II） & \hyperref[ex:6.2.27]{235} \\
例 6.3.2（负二项似然） & \hyperref[ex:6.3.2]{236} \\
例 6.3.3（正态信仰分布） & \hyperref[ex:6.3.3]{237} \\
例 6.3.4（证据函数） & \hyperref[ex:6.3.4]{238} \\
例 6.3.5（二项／负二项试验） & \hyperref[ex:6.3.5]{238} \\
例 6.3.7（例6.3.5续） & \hyperref[ex:6.3.7]{240} \\
例 6.4.1（二项等变性） & \hyperref[ex:6.4.1]{242} \\
例 6.4.3（例6.4.1续） & \hyperref[ex:6.4.3]{242} \\
例 6.4.5（例6.4.1的结论） & \hyperref[ex:6.4.5]{243} \\
例 6.4.6（正态位置不变性） & \hyperref[ex:6.4.6]{243} \\
\addlinespace[.5em]
例 7.2.1（正态分布的矩方法） & \hyperref[ex:7.2.1]{253} \\
例 7.2.2（二项分布的矩方法） & \hyperref[ex:7.2.2]{254} \\
例 7.2.3（Satterthwaite近似） & \hyperref[ex:7.2.3]{254} \\
例 7.2.5（正态似然） & \hyperref[ex:7.2.5]{256} \\
例 7.2.6（例7.2.5续） & \hyperref[ex:7.2.6]{257} \\
例 7.2.7（Bernoulli最大似然估计量） & \hyperref[ex:7.2.7]{257} \\
例 7.2.8（取值范围受限的最大似然估计量） & \hyperref[ex:7.2.8]{257} \\
例 7.2.9（试验次数未知的二项最大似然估计量） & \hyperref[ex:7.2.9]{258} \\
例 7.2.11（$\mu$与$\sigma$未知时的正态最大似然估计量） & \hyperref[ex:7.2.11]{260} \\
例 7.2.12（例7.2.11续） & \hyperref[ex:7.2.12]{260} \\
例 7.2.13（例7.2.2续） & \hyperref[ex:7.2.13]{261} \\
例 7.2.14（二项Bayes估计） & \hyperref[ex:7.2.14]{262} \\
例 7.2.16（正态Bayes估计量） & \hyperref[ex:7.2.16]{263} \\
例 7.2.17（多个Poisson发生率） & \hyperref[ex:7.2.17]{264} \\
例 7.2.18（例7.2.17续） & \hyperref[ex:7.2.18]{265} \\
例 7.2.19（例7.2.17的结论） & \hyperref[ex:7.2.19]{266} \\
例 7.3.3（正态均方误差） & \hyperref[ex:7.3.3]{267} \\
例 7.3.4（例7.3.3续） & \hyperref[ex:7.3.4]{268} \\
例 7.3.5（二项Bayes估计量的均方误差） & \hyperref[ex:7.3.5]{269} \\
例 7.3.6（等变估计量的均方误差） & \hyperref[ex:7.3.6]{270} \\
例 7.3.8（Poisson无偏估计） & \hyperref[ex:7.3.8]{271} \\
例 7.3.12（例7.3.8的结论） & \hyperref[ex:7.3.12]{274} \\
例 7.3.13（尺度均匀分布的无偏估计量） & \hyperref[ex:7.3.13]{275} \\
例 7.3.14（正态方差界） & \hyperref[ex:7.3.14]{276} \\
例 7.3.16（例7.3.14续） & \hyperref[ex:7.3.16]{276} \\
例 7.3.18（以非充分统计量为条件） & \hyperref[ex:7.3.18]{278} \\
例 7.3.21（零的无偏估计量） & \hyperref[ex:7.3.21]{279} \\
例 7.3.22（例7.3.13续） & \hyperref[ex:7.3.22]{280} \\
例 7.3.24（二项最佳无偏估计） & \hyperref[ex:7.3.24]{281} \\
例 7.3.25（二项风险函数） & \hyperref[ex:7.3.25]{283} \\
例 7.3.26（正态方差的风险） & \hyperref[ex:7.3.26]{283} \\
例 7.3.27（使用Stein损失的方差估计） & \hyperref[ex:7.3.27]{284} \\
例 7.3.28（两个Bayes规则） & \hyperref[ex:7.3.28]{285} \\
例 7.3.29（正态Bayes估计） & \hyperref[ex:7.3.29]{286} \\
例 7.3.30（二项Bayes估计） & \hyperref[ex:7.3.30]{286} \\
\addlinespace[.5em]
例 8.2.2（正态似然比检验） & \hyperref[ex:8.2.2]{304} \\
例 8.2.3（指数分布的似然比检验） & \hyperref[ex:8.2.3]{304} \\
例 8.2.5（似然比检验与充分性） & \hyperref[ex:8.2.5]{306} \\
例 8.2.6（方差未知的正态似然比检验） & \hyperref[ex:8.2.6]{306} \\
例 8.2.7（正态Bayes检验） & \hyperref[ex:8.2.7]{307} \\
例 8.2.8（正态并交检验） & \hyperref[ex:8.2.8]{308} \\
例 8.2.9（验收抽样） & \hyperref[ex:8.2.9]{309} \\
例 8.3.2（二项功效函数） & \hyperref[ex:8.3.2]{310} \\
例 8.3.3（正态功效函数） & \hyperref[ex:8.3.3]{311} \\
例 8.3.4（例8.3.3续） & \hyperref[ex:8.3.4]{312} \\
例 8.3.7（似然比检验的大小） & \hyperref[ex:8.3.7]{313} \\
例 8.3.8（并交检验的大小） & \hyperref[ex:8.3.8]{313} \\
例 8.3.10（例8.3.3的结论） & \hyperref[ex:8.3.10]{314} \\
例 8.3.14（二项一致最强检验） & \hyperref[ex:8.3.14]{316} \\
例 8.3.15（正态一致最强检验） & \hyperref[ex:8.3.15]{317} \\
例 8.3.18（例8.3.15续） & \hyperref[ex:8.3.18]{318} \\
例 8.3.19（一致最强检验不存在） & \hyperref[ex:8.3.19]{318} \\
例 8.3.20（无偏检验） & \hyperref[ex:8.3.20]{319} \\
例 8.3.22（一个等价关系） & \hyperref[ex:8.3.22]{320} \\
例 8.3.25（交并检验） & \hyperref[ex:8.3.25]{321} \\
例 8.3.28（双侧正态$p$值） & \hyperref[ex:8.3.28]{322} \\
例 8.3.29（单侧正态$p$值） & \hyperref[ex:8.3.29]{323} \\
例 8.3.30（Fisher精确检验） & \hyperref[ex:8.3.30]{323} \\
例 8.3.31（一致最强检验的风险） & \hyperref[ex:8.3.31]{325} \\
\addlinespace[.5em]
例 9.1.2（区间估计量） & \hyperref[ex:9.1.2]{337} \\
例 9.1.3（例9.1.2续） & \hyperref[ex:9.1.3]{338} \\
例 9.1.6（尺度均匀分布的区间估计量） & \hyperref[ex:9.1.6]{338} \\
例 9.2.1（反演正态检验） & \hyperref[ex:9.2.1]{339} \\
例 9.2.3（反演似然比检验） & \hyperref[ex:9.2.3]{342} \\
例 9.2.4（正态单侧置信界） & \hyperref[ex:9.2.4]{343} \\
例 9.2.5（二项单侧置信界） & \hyperref[ex:9.2.5]{344} \\
例 9.2.7（位置–尺度枢轴量） & \hyperref[ex:9.2.7]{345} \\
例 9.2.8（伽马枢轴量） & \hyperref[ex:9.2.8]{346} \\
例 9.2.9（例9.2.8续） & \hyperref[ex:9.2.9]{347} \\
例 9.2.10（正态枢轴区间） & \hyperref[ex:9.2.10]{347} \\
例 9.2.11（最短长度二项集合） & \hyperref[ex:9.2.11]{348} \\
例 9.2.13（位置指数分布的区间） & \hyperref[ex:9.2.13]{350} \\
例 9.2.15（Poisson区间估计量） & \hyperref[ex:9.2.15]{351} \\
例 9.2.16（Poisson可信集） & \hyperref[ex:9.2.16]{352} \\
例 9.2.17（Poisson可信概率与覆盖概率） & \hyperref[ex:9.2.17]{354} \\
例 9.2.18（正态可信集的覆盖概率） & \hyperref[ex:9.2.18]{354} \\
例 9.3.1（使长度最优） & \hyperref[ex:9.3.1]{356} \\
例 9.3.3（使期望长度最优） & \hyperref[ex:9.3.3]{358} \\
例 9.3.4（最短枢轴区间） & \hyperref[ex:9.3.4]{359} \\
例 9.3.6（一致最准确置信界） & \hyperref[ex:9.3.6]{360} \\
例 9.3.8（例9.3.6续） & \hyperref[ex:9.3.8]{361} \\
例 9.3.11（Poisson最高后验密度区域） & \hyperref[ex:9.3.11]{362} \\
例 9.3.12（正态最高后验密度区域） & \hyperref[ex:9.3.12]{363} \\
例 9.3.13（正态区间估计量） & \hyperref[ex:9.3.13]{364} \\
\addlinespace[.5em]
例 10.1.2（$\bar X$的相合性） & \hyperref[ex:10.1.2]{378} \\
例 10.1.4（例10.1.2续） & \hyperref[ex:10.1.4]{379} \\
例 10.1.8（极限方差） & \hyperref[ex:10.1.8]{379} \\
例 10.1.10（大样本混合方差） & \hyperref[ex:10.1.10]{380} \\
例 10.1.13（渐近正态性与相合性） & \hyperref[ex:10.1.13]{381} \\
例 10.1.14（近似二项方差） & \hyperref[ex:10.1.14]{382} \\
例 10.1.15（例10.1.14续） & \hyperref[ex:10.1.15]{383} \\
例 10.1.17（Poisson估计量的渐近相对效率） & \hyperref[ex:10.1.17]{384} \\
例 10.1.18（估计伽马分布的均值） & \hyperref[ex:10.1.18]{385} \\
例 10.1.19（用自助法估计方差） & \hyperref[ex:10.1.19]{386} \\
例 10.1.20（用自助法估计二项方差） & \hyperref[ex:10.1.20]{387} \\
例 10.1.21（例10.1.20的结论） & \hyperref[ex:10.1.21]{387} \\
例 10.1.22（参数自助法） & \hyperref[ex:10.1.22]{388} \\
例 10.2.1（样本均值的稳健性） & \hyperref[ex:10.2.1]{389} \\
例 10.2.3（中位数的渐近正态性） & \hyperref[ex:10.2.3]{390} \\
例 10.2.4（中位数相对于均值的渐近相对效率） & \hyperref[ex:10.2.4]{391} \\
例 10.2.5（Huber估计量） & \hyperref[ex:10.2.5]{391} \\
例 10.2.6（Huber估计量的极限分布） & \hyperref[ex:10.2.6]{392} \\
例 10.2.7（Huber估计量的渐近相对效率） & \hyperref[ex:10.2.7]{393} \\
例 10.3.2（Poisson似然比检验） & \hyperref[ex:10.3.2]{395} \\
例 10.3.4（多项似然比检验） & \hyperref[ex:10.3.4]{396} \\
例 10.3.5（大样本二项检验） & \hyperref[ex:10.3.5]{398} \\
例 10.3.6（二项得分检验） & \hyperref[ex:10.3.6]{400} \\
例 10.3.7（基于Huber估计量的检验） & \hyperref[ex:10.3.7]{400} \\
例 10.4.1（例10.1.14续） & \hyperref[ex:10.4.1]{402} \\
例 10.4.2（二项得分区间） & \hyperref[ex:10.4.2]{403} \\
例 10.4.3（二项似然比检验区间） & \hyperref[ex:10.4.3]{403} \\
例 10.4.4（近似区间） & \hyperref[ex:10.4.4]{404} \\
例 10.4.5（近似Poisson区间） & \hyperref[ex:10.4.5]{404} \\
例 10.4.6（二项得分区间的进一步讨论） & \hyperref[ex:10.4.6]{405} \\
例 10.4.7（二项区间的比较） & \hyperref[ex:10.4.7]{406} \\
例 10.4.8（基于Huber估计量的区间） & \hyperref[ex:10.4.8]{406} \\
例 10.4.9（负二项区间） & \hyperref[ex:10.4.9]{408} \\
例 10.6.2（均值与中位数的影响函数） & \hyperref[ex:10.6.2]{419} \\
\addlinespace[.5em]
例 11.2.1（单因素方差分析） & \hyperref[ex:11.2.1]{422} \\
例 11.2.3（方差分析假设） & \hyperref[ex:11.2.3]{424} \\
例 11.2.6（方差分析中的对比） & \hyperref[ex:11.2.6]{427} \\
例 11.2.9（成对差） & \hyperref[ex:11.2.9]{431} \\
例 11.2.12（例11.2.1续） & \hyperref[ex:11.2.12]{434} \\
例 11.3.1（预测葡萄产量） & \hyperref[ex:11.3.1]{436} \\
例 11.3.4（例11.3.1续） & \hyperref[ex:11.3.4]{449} \\
\addlinespace[.5em]
例 12.2.1（估计大气压） & \hyperref[ex:12.2.1]{468} \\
例 12.3.1（挑战者号数据） & \hyperref[ex:12.3.1]{481} \\
例 12.3.2（挑战者号数据续） & \hyperref[ex:12.3.2]{482} \\
例 12.4.1（最小二乘估计的稳健性） & \hyperref[ex:12.4.1]{484} \\
例 12.4.2（灾难性观测） & \hyperref[ex:12.4.2]{484} \\
例 12.4.3（最小绝对偏差估计量的渐近正态性） & \hyperref[ex:12.4.3]{484} \\
例 12.4.4（回归M估计量） & \hyperref[ex:12.4.4]{487} \\
例 12.4.5（回归渐近相对效率的模拟） & \hyperref[ex:12.4.5]{487} \\
例 12.6.1（无序抽样） & \hyperref[ex:12.6.1]{497} \\
例 12.6.2（一元变换） & \hyperref[ex:12.6.2]{498} \\
例 12.6.3（二元变换） & \hyperref[ex:12.6.3]{498} \\
例 12.6.4（正态概率） & \hyperref[ex:12.6.4]{499} \\
例 12.6.5（和的密度） & \hyperref[ex:12.6.5]{500} \\
例 12.6.6（均匀随机变量之和的四阶矩） & \hyperref[ex:12.6.6]{501} \\
例 12.6.7（伽马均值的渐近相对效率） & \hyperref[ex:12.6.7]{501} \\
例 12.6.8（卡方矩母函数的极限） & \hyperref[ex:12.6.8]{502} \\
\end{longtable}

% 用户授权疑误脚注：2026-11-04
